\section{Integer and fixed-point formats}
\label{sec:formats}

The float32 exploration ended with the DSPs nearly used up: each float multiply-add takes five DSP48E1 slices
(\cref{sec:eval}). Narrower element formats need fewer of them, so the same device can hold far larger arrays. This
section describes how the flow was extended to integer and fixed-point formats \emph{without rewriting the kernel},
and what twelve more campaigns measured on the card (24 November 2026, \cref{tab:fmt-campaigns}). The dated record
is in \file{docs/notes2/}.

\subsection{One element type, no kernel change}
\label{sec:formats:approach}
\code{gemm\_hls} takes its element type from one line of \file{Config.h} (\code{using Data\_t = ...}), and uses it for
A, B, C \emph{and the accumulation}: the product and the running sum have the same type. An integer GEMM therefore
wraps at the element width, as C integer arithmetic does. That makes the results exactly checkable, and it measures
what the hardware costs and how fast it runs. A quantised network would also need a wider accumulator, which
is a kernel change (\cref{sec:lessons}). \Cref{tab:tiers} lists what works through the generated \file{Config.h} alone.

\begin{table}[ht]
\centering\small
\caption{Element formats by how much they need. \(W\): bits per element, \(F\): fractional bits. Elements per
  64-byte bus word follow from \(W\), since every m\_axi port of the kernel must stay 512 bits wide.}
\label{tab:tiers}
\begin{tabular}{@{}llllp{5.2cm}@{}}
\toprule
Tier & Formats & C type & Per word & Status \\
\midrule
T0 & float32 & \code{float} & 16 & \cref{sec:eval} \\
T1 & int/uint, \(W\) = 8, 16, 32, 64 & \code{int8\_t} \dots{} \code{uint64\_t} & 64 \dots{} 8 & unchanged kernel; int64 synthesised only \\
T2 & fixed/ufixed, \(W\) = 16, 32 & \code{ap\_fixed<W,W-F>} & 32, 16 & \code{\#include <ap\_fixed.h>} added to the generated header; 8-bit fixed point crashes Vitis HLS 2026.1 \\
T3 & int/uint, \(W\) = 2, 4 & \code{ap\_int<W>} & 256, 128 & packed lanes; the kernel is given a \emph{logical} bus width of \(512/W\) bytes \\
\midrule
-- & \(W\) = 12, 24, \dots & & & port width \(\ne\) 512: the kernel would have to count elements by bits \\
-- & wider accumulator, pow2 & & & separate A/B/C types and a shift operator (DenseShift~\cite{denseshift}) \\
\bottomrule
\end{tabular}
\end{table}

Tier T3 rests on a detail of hlslib: \code{DataPack} packs \code{ap\_int} elements densely by their exact bit width,
while the kernel counts elements per bus word from \code{sizeof(Data\_t)}, which is one byte for a 4-bit type.
Telling the kernel that the bus is \(512/W\) bytes wide makes the two agree: 128 int4 elements per 512-bit word, which
HLS confirmed for every 4-bit point.

\subsection{Flow and host support}
\begin{itemize}
  \item \textbf{Configuration.} \code{kernel.data\_type} is a format family (\code{float}, \code{int}, \code{uint},
        \code{fixed}, \code{ufixed}) with \code{data\_bits} and \code{frac\_bits}; \code{accum\_bits} is reserved.
        The new fields enter the build hash only when they are set, so all 24 existing configurations kept their build
        ids, HLS hashes and a byte-identical \file{Config.h}: no float build was invalidated. Validation, size multiples
        and \file{Config.h} widths all work in elements, and unsupported formats are rejected with the reason.
  \item \textbf{Checks in the flow.} The HLS stage fails a point whose m\_axi ports are not 512 bits wide, and a
        resource gate records a point as \code{over\_budget} without a Vivado run when the HLS estimate needs more
        than 1910 BRAM18 (with 150 for the SoC) or 1920 DSPs. Linked grid axes (\code{"@format"}, \code{"@array"})
        keep format and array sweeps free of meaningless combinations.
  \item \textbf{Host.} \code{gemm\_runner --dtype int4} (also \code{fixed16.8}, \code{uint8}, \dots) draws data over the
        full code range, packs it in \code{DataPack}'s bit order, and checks every result exactly against a model of the
        kernel's arithmetic: products and sums wrap at \(W\) bits; fixed-point products are floored by \(2^F\)
        (\code{AP\_TRN}, \code{AP\_WRAP}). The model matched Vitis' own \code{ap\_int}/\code{ap\_fixed} on 2000 random
        16-term dot products in each of 11 formats.
  \item \textbf{Results.} Every build and run row carries the format (\code{dtype}, kind, bits, signedness,
        fractional bits, elements per word, and A/B/C and accumulator types, which are all equal for now) and
        \code{gops}/\code{peak\_gops}.
\end{itemize}

\subsection{Cost of a format at equal throughput}
The HLS probe (70 valid points, PE 16\x{}8 to 16\x{}128) showed every format with 512-bit ports, and gave two negative
results. \code{mult\_resource: fabric} (\code{config\_op mul -impl fabric}) also moves the kernel's 32-bit address
multiplies into LUTs, and the clock estimate jumps to 7.6\,ns, so it was dropped. 8-bit \code{ap\_fixed} makes Vitis
HLS crash, so the flow rejects it.

\Cref{tab:same,fig:formats_cost} compare the formats on the array of the best float build, PE 32\x8 with tile
1024\x512 at 200\,MHz. The throughput is the same for every format: the array, not the format, sets the rate. The cost
falls sharply with width. Vivado maps the 8-bit multiplies to LUTs, so int8 needs 40 DSPs (address logic only)
against float's 1320, and under half of float's LUTs and BRAM. int16 needs one DSP per MAC. fixed32.16 needs 1064 DSPs
against int32's 808, because the full 64-bit product is formed before the shift by 16.

\begin{table}[ht]
\centering\small
\caption{One array in five formats: PE 32\x8, tile 1024\x512, transposed A, 200\,MHz, default preset
  (\file{data/formats\_same.csv}). Peak 102.4\,GOP/s for all.}
\label{tab:same}
\begin{tabular}{@{}lrrrrrrl@{}}
\toprule
Format & DSP & LUT & FF & BRAM tiles & kernel LUT & GOP/s at \(8192^3\) & Verification \\
\midrule
float32 & 1320 & 154.3\,k & 225.4\,k & 654.5 & 125.6\,k & 98.59 & exact after DDR tuning \\
int8 & \textbf{40} & 69.4\,k & 75.7\,k & 299.5 & 40.8\,k & 98.64 & exact \\
int16 & 296 & 67.0\,k & 83.0\,k & 428.0 & 38.3\,k & 98.64 & exact \\
fixed16.8 & 296 & 67.3\,k & 79.0\,k & 428.5 & 38.6\,k & 98.64 & exact \\
int32 & 808 & 81.4\,k & 102.2\,k & 654.5 & 52.6\,k & 98.59 & exact after DDR tuning \\
\bottomrule
\end{tabular}
\end{table}

\begin{figure}[ht]
  \centering
  % Resources of PE 32x8, tile 1024x512, 200 MHz in each element format, in % of the float32 build.
\begin{tikzpicture}
\begin{axis}[rbar, ybar=1pt, bar width=6.5pt, height=5.6cm, ymax=115, ylabel={\% of the float32 build},
    xtick=data, xticklabels={float32, int8, int16, fixed16.8, int32}, enlarge x limits=0.13,
    legend columns=4, legend style={at={(0.5,1.02)}, anchor=south}]
  \addplot[s1] table[x expr=\coordindex, y=dsp_rel]  {data/formats_same.csv}; \addlegendentry{DSP}
  \addplot[s2] table[x expr=\coordindex, y=lut_rel]  {data/formats_same.csv}; \addlegendentry{LUT}
  \addplot[s3] table[x expr=\coordindex, y=ff_rel]   {data/formats_same.csv}; \addlegendentry{FF}
  \addplot[s5] table[x expr=\coordindex, y=bram_rel] {data/formats_same.csv}; \addlegendentry{BRAM tiles}
\end{axis}
\end{tikzpicture}

  \caption{Resources of the same array (PE 32\x8, tile 1024\x512, 200\,MHz) in each format, in \% of the float32
    build.}
  \label{fig:formats_cost}
\end{figure}

\subsection{Wider arrays}
\label{sec:formats:wide}
With the DSPs freed, the arrays could grow until BRAM, LUTs or timing stopped them (\cref{tab:formats,fig:formats_best}):
\begin{itemize}
  \item \textbf{int8:} PE 32\x32 (1024 MACs) reaches \SI{379.4}{GOP/s} at 200\,MHz and \textbf{\SI{426.6}{GOP/s} at
        225\,MHz}, 3.8 times the best float result, with 40 DSPs and \SI{44}{\percent} of the LUTs.
  \item \textbf{int16:} PE 32\x16 (\SI{190.9}{GOP/s}). PE 32\x32 needs 2045 BRAM18 by the HLS estimate, because the
        C tile is twice as wide per element as int8's. The resource gate stopped it before Vivado.
  \item \textbf{int4:} PE 32\x64 (2048 MACs) reaches \SI{756.7}{GOP/s} at 200\,MHz and \SI{854.1}{GOP/s} at 225\,MHz.
        PE 32\x128 is over the BRAM budget (2045 BRAM18). \textbf{PE 64\x64 (4096 MACs) reaches \SI{1290.1}{GOP/s}
        at 180\,MHz}, 11.6 times the best float result, with 38 DSPs, \SI{67}{\percent} of the LUTs and
        \SI{68}{\percent} of the BRAM. HLS overestimates LUTs for the 4-bit arrays by 2--3\x{} (PE 32\x64: 240\,k
        estimated, 117\,k placed).
\end{itemize}

\paragraph{Timing.} With 32 PEs, int8 and int4 both close 225\,MHz with the default preset; int8 misses 250\,MHz
by 0.365\,ns. The limiting path is no longer in the PEs. It is the control of the HLS dataflow region:
a PE's loop counter drives the reset of the \code{ap\_sync\_reg\_*\_ap\_ready} flops of every process in the region
(6 logic levels, \SI{81}{\percent} routing), and that fan-out grows with the number of PEs. The three PE 64\x64
builds at 200\,MHz missed by 0.215, 0.051 and 0.336\,ns on this path. At 180\,MHz the default preset met the accel
clock but failed the PCIe \SI{250}{MHz} user clock by 0.447\,ns: the path from a LitePCIe FIFO in block RAM into the
PCIe block's transmit throttle, displaced by the congested design. The \code{aggressive} and \code{explore} presets
close both, with +0.045 and +0.018\,ns. Going faster with 64 PEs needs a kernel or directive change in the
dataflow control (\cref{sec:lessons}).

\begin{table}[ht]
\centering\small
\caption{Best verified build per format (kernel time). Peak = \(2 P_N P_M f\). Every run is checked exactly
  (float32: bit for bit). The int32, fixed32.16 and both int4 builds also passed a full screen (\(2\times2048^3\), every
  entry), the int8 build a full \(1024^3\) check.}
\label{tab:formats}
\begin{tabular}{@{}lllrrrl@{}}
\toprule
Format & Array, tile & Clock, preset & Peak & Best (size) & \% & LUT / DSP / BRAM tiles \\
\midrule
float32 & PE 32\x8, 1024\x512 & 225\,MHz, aggr. & 115.2 & 110.9 (\(8192^3\)) & 96.2 & 152.0\,k / 1320 / 654.5 \\
fixed32.16 & PE 32\x8, 1024\x512 & 200\,MHz, aggr. & 102.4 & 98.6 (\(8192^3\)) & 96.3 & 79.2\,k / 1064 / 654.5 \\
int32 & PE 32\x8, 1024\x512 & 200\,MHz, aggr. & 102.4 & 98.2 (\(4096^3\)) & 95.9 & 81.3\,k / 808 / 654.5 \\
int16 & PE 32\x16, 512\x1024 & 200\,MHz & 204.8 & 190.9 (\(8192^3\)) & 93.2 & 85.7\,k / 552 / 414.5 \\
int8 & PE 32\x32, 512\x1024 & 225\,MHz & 460.8 & 426.6 (\(8192^3\)) & 92.6 & 131.3\,k / 40 / 414.5 \\
int4 & PE 32\x64, 512\x2048 & 225\,MHz & 921.6 & 854.1 (\(8192^3\)) & 92.7 & 117.5\,k / 39 / 414.5 \\
\textbf{int4} & \textbf{PE 64\x64, 512\x4096} & \textbf{180\,MHz, aggr.} & 1474.6 & \textbf{1290.1} (\(8192^3\)) & 87.5 & 200.3\,k / 38 / 652.5 \\
\bottomrule
\end{tabular}
\end{table}

\begin{figure}[ht]
  \centering
  % Best verified GOP/s per element format, against the peak 2*PN*PM*f of each build.
\begin{tikzpicture}
\begin{axis}[rpt, ybar=0pt, bar width=16pt, height=6.4cm, ymin=0, ymax=1650,
    ylabel={GOP/s (kernel)}, xtick=data, enlarge x limits=0.09,
    x tick label style={align=center, font=\sffamily\tiny},
    xticklabels={{float32\\PE 32\x8, 225\,MHz\\1320 DSP}, {fixed32.16\\PE 32\x8, 200\\1064 DSP},
                 {int32\\PE 32\x8, 200\\808 DSP}, {int16\\PE 32\x16, 200\\552 DSP},
                 {int8\\PE 32\x32, 225\\40 DSP}, {int4\\PE 32\x64, 225\\40 DSP}, {int4\\PE 64\x64, 180\\38 DSP}},
    nodes near coords, nodes near coords style={font=\sffamily\scriptsize, /pgf/number format/.cd, fixed, precision=0},
    legend style={at={(0.02,0.98)}, anchor=north west}]
  \addplot[s7, fill=black!8, draw=black!45, bar shift=0pt, nodes near coords=\empty] table[x=i, y=peak] {data/formats_best.csv};
  \addlegendentry{peak of the build}
  \addplot[s1, bar shift=0pt, bar width=10pt] table[x=i, y=gops] {data/formats_best.csv};
  \addlegendentry{best measured}
\end{axis}
\end{tikzpicture}

  \caption{Best kernel GOP/s per format (\file{data/formats\_best.csv}) against the peak of each build. The DSP count
    of each build is under its label.}
  \label{fig:formats_best}
\end{figure}

\subsection{Wrong results under load: DDR3 read margin}
\label{sec:formats:ddr}
Exact integer checks with full-range data exposed errors that the float tolerance had hidden. The int32 build of
\cref{tab:same} returned 10--\SI{19}{\percent} wrong entries and int8 PE 16\x64 4--\SI{6}{\percent}, both with timing met.
\begin{itemize}
  \item \textbf{Symptoms.} Whole C rows (int32) or columns (int8 PE 16\x64) inside a tile were wrong. With B = I or
        A = I, the results were exact. Errors needed both operands switching, which first suggested a timing failure
        under heavy activity. But a lower clock (180\,MHz) did not help.
  \item \textbf{Which bits.} Recovering the corrupted inputs from the dumped matrices showed single-bit flips at
        word bit \(\equiv 5 \pmod{64}\): DDR3 data line DQ5 of channel 0, on the odd beats of a burst. The int8 and
        float errors map to the same kind of lines and beats.
  \item \textbf{Cause.} The DDR3 PHY samples the data lines with an internal clock (not DQS), with one delay tap per
        byte lane. Calibration at bring-up finds a window of about 11 taps on an idle chip; under GEMM load it shrinks
        to about 3, and the calibrated tap sits at its edge. Moving only that lane's tap changed the error rate from
        \(\sim\)17\,k to 524\,k of 1\,M entries. DMA streaming at the same taps with an idle core is clean, so it is
        the core's switching activity, not the traffic, that shifts the capture point.
  \item \textbf{Fix (host only).} \code{gemmflow.ddr\_tune} keeps the bring-up calibration, then sweeps each lane's tap
        around it while the build's own kernel runs an exact check, and takes the centre of the error-free run
        (\code{gemm\_run.py --ddr-load-tune}, campaign key \code{"ddr\_load\_tune"}). All four failing builds became
        exact (0 bad over \(2 \times 2048^3\), 8.4\,M entries each), with unchanged throughput. Every campaign started
        after the fix ran with it. A permanent fix belongs in the PHY (per-bit delays, DQS capture) or the calibration (under load),
        and would rebuild every bitstream.
\end{itemize}
Bit-exact screens (\code{--float-exact}, the host accumulating in the kernel's order) of all float builds then found
one affected build: PE 32\x8 at 200\,MHz (\code{a88777ec}), with errors of \(\sim\)\(4\cdot10^{-6}\) relative,
within the tolerance of \cref{sec:host}. Its throughput is unaffected, and it is exact after tuning. The headline
225\,MHz build and 17 others are bit-exact without tuning. Float runs are now checked bit for bit.

\subsection{Two tile rules the kernel does not check}
\label{sec:formats:tiles}
An int4 PE 48\x64 build with tile 768\x2048 met timing, returned \SI{94}{\percent} wrong entries on its first check,
and then hung on every run. The cause is two assumptions in \code{gemm\_hls}'s \file{Compute.cpp} that neither the
upstream build nor our flow checked:
\begin{enumerate}
  \item \textbf{\(T_M/P_M \ge P_N\).} During each outer product a PE passes the next A values down the chain in the slots
        \(m_1 = \text{its index} \dots P_N-1\) of a loop that is only \(T_M/P_M\) long. If that is shorter than
        \(P_N\), the PEs at the end never get their values: C is wrong, and the values left over block the feeder.
        PE 48\x64 had \(2048/64 = 32 < 48\).
  \item \textbf{\(T_M/P_M\) and \(T_N/P_N\) are powers of two.} The C write-back counts in
        \code{ap\_uint<ConstLog2(x)>}, which rounds down.
\end{enumerate}
A g++ simulation of the upstream testbench on ten small arrays was wrong exactly when a rule was broken, including
a non-power-of-two \(P_N\) = 12, which works when both rules hold. Both rules are now validation errors.
Every earlier board build met both. PE 64\x64 then used tile 512\x4096 (\(T_M/P_M = 64 = P_N\)), the build in
\cref{tab:formats}.

\subsection{A fixed loss per k step}
\label{sec:formats:kstep}
PE 64\x64 reaches only \SI{87.5}{\percent} of its peak, against \SI{92}{\percent} for the 32-PE arrays, and the gap
does not shrink from \(4096^3\) to \(8192^3\). So it is a cost per step of the k loop, not per tile. From the measured
cycle counts, each build loses a roughly fixed number of cycles per k step and memory tile
(\cref{tab:kstep}), where the ideal is \((T_N/P_N)(T_M/P_M)\) cycles. The float PE 32\x8 already showed this: doubling
its tile raised it from 93.2 to \SI{96.3}{\percent} with about the same loss per step. Where those cycles go in the
kernel is not yet known: pipeline fill of the PE loop or a restart of the A/B feed. Tiles with twice the work per
k step are the next experiment for the int4 arrays.

\begin{table}[ht]
\centering\small
\caption{Cycles lost per k step and memory tile, at the largest measured size.}
\label{tab:kstep}
\begin{tabular}{@{}lrrr@{}}
\toprule
Build & Ideal cycles per k step & Lost per k step & \% of peak \\
\midrule
float32 PE 32\x8, tile 512\x512 & 1024 & 75 & 93.2 \\
float32 / int8 PE 32\x8, tile 1024\x512 & 2048 & 79 & 96.3 \\
int8 PE 32\x32, tile 512\x1024 & 512 & 41 & 92.6 \\
int4 PE 32\x64, tile 512\x2048 & 512 & 41 & 92.7 \\
int4 PE 64\x64, tile 512\x4096 & 512 & 73 & 87.5 \\
\bottomrule
\end{tabular}
\end{table}

\paragraph{Sizes and transfers.} A size sweep on the int8 PE 32\x8 build showed the same whole-tile behaviour as float
(\cref{sec:eval:sizes}). Narrow formats move fewer bytes: at \(1024^3\) the host spends 5.6\,ms on transfers and
packing (\SI{20}{\percent} of the end-to-end time), against 15.5\,ms (\SI{41}{\percent}) for float32 on the same array.

\subsection{Campaigns}
\begin{table}[ht]
\centering\small
\caption{Integer and fixed-point campaigns, 24 November 2026 (\file{data/campaigns.csv}). Builds: new Vivado runs;
  runs: result rows (15 of them with the DDR read errors of \cref{sec:formats:ddr}).}
\label{tab:fmt-campaigns}
\begin{tabular}{@{}lllrrrl@{}}
\toprule
Campaign & Time & Stages & Points & Builds & Runs & Question \\
\midrule
\code{dtype\_hls} & 03:25--04:53 & sim, hls & 70 & -- & -- & which formats synthesise, at what cost? \\
\code{dtype\_same\_array} & 04:53--05:53 & all & 5 & 4 & 30 & each format on the same array \\
\code{dtype\_size\_sweep} & 05:53 & run & 1 & -- & 90 & 45 matrix shapes in int8 \\
\code{int8\_scale} & 05:53--06:42 & all & 3 & 3 & 18 & widest int8 array \\
\code{int16\_scale} & 06:42--07:13 & all & 3 & 2 & 12 & widest int16 array \\
\code{int32\_margin} & 07:13--07:47 & all & 2 & 2 & 8 & int32 errors: timing margin? \\
\code{int4\_probe} & 07:46--09:07 & all & 4 & 4 & 16 & packed 4-bit lanes end to end \\
\code{fixed\_point} & 09:07--11:38 & all & 2 & 2 & 6 & fixed32.16 \\
\code{dtype\_clock} & 09:26--11:20 & all & 6 & 5 & 12 & int8 PE 32\x32 at 200/225/250\,MHz \\
\code{int4\_wide} & 11:38--13:47 & all & 3 & 2 & 0 & 4096-MAC int4 arrays (stopped, \cref{sec:formats:tiles}) \\
\code{int4\_pe64} & 15:33--21:12 & all & 5 & 5 & 8 & PE 64\x64 with a legal tile \\
\code{int4\_clock} & 18:36--21:48 & all & 2 & 2 & 8 & int4 PE 32\x64 at 225\,MHz \\
\midrule
\multicolumn{3}{@{}l}{total} & & 31 & 208 & \\
\bottomrule
\end{tabular}
\end{table}
